% F. Resultados obtenidos, justificando los no alcanzados (Guía apdo. 3.F)
% Extensión orientativa: 5-7 páginas.
% REGLA: aquí solo van resultados MEDIDOS. Toda cifra debe proceder de una campaña
% registrada en docs/EXPERIMENTOS.md. Nunca estimaciones ni valores plausibles.
\chapter{Resultados}
\label{cap:resultados}



\section{Caracterización de la cadena de adquisición}

Antes de exponer los resultados sobre el activo procede consignar los de la cadena de medida,
dado que acotan lo que las demás cifras pueden significar. Se detallan en el
anexo~\ref{anexo:verificacion} y aquí se retienen los dos que condicionan la interpretación.

El primero es que la frecuencia fundamental del activo se identifica con un coeficiente de
variación del \mbox{0,15 \%}, lo que constituye por sí mismo una capacidad de medida: permite
seguir variaciones pequeñas del régimen de giro. La repetibilidad en amplitud, del \mbox{7,4 \%},
sitúa en cambio el umbral de detección en un cambio del orden del \mbox{21 \%} para esa magnitud.

El segundo es que se localizaron dos defectos que corrompían el dato \textbf{en silencio}, ambos
por bibliotecas que no propagan los errores del bus de comunicación. Una sola muestra corrupta
entre las 1024 de una ráfaga elevaba la kurtosis de 2,9 a 847, valor que coincide con la
predicción analítica para un único impulso aislado. La gravedad no reside en la pérdida de una
medida sino en que \textbf{el artefacto imita el fenómeno que el sistema debe detectar}: la
kurtosis se seleccionó precisamente por su sensibilidad a los impulsos aislados.

\section{Detección de un fallo real}

Durante la fase de captura se dispuso de un segundo activo del mismo tipo que presentaba un
fallo real, no inducido, perceptible como un tono audible sostenido. La circunstancia permitió
contrastar el sistema frente a un fallo con referencia de verdad independiente de la propia
instrumentación, y frente a un activo de referencia en estado nominal medido con el mismo
firmware en el mismo intervalo temporal. Los resultados proceden de las campañas EXP-003 y
EXP-004.

El cuadro~\ref{table:firma-fallo} recoge la firma detectada.

\begin{table}[ht!]
\caption{Firma del fallo real (campaña EXP-003, 7,14 h, 815 ráfagas, 676 utilizables).}
\label{table:firma-fallo}
\centering
\begin{tabular}{|p{6.4cm}|p{6.0cm}|}
\hline
\textbf{Magnitud} & \textbf{Valor} \\
\hline
Frecuencia de giro & \mbox{49,745 Hz} (CV 0,12 \%) \\
\hline
Componentes tonales & \mbox{398 Hz}, \mbox{448 Hz} y \mbox{497 Hz} \\
\hline
Relación con la frecuencia de giro & \textbf{8}, \textbf{9} y \textbf{10} \\
\hline
Desviación respecto del entero & inferior al 0,05 \% \\
\hline
Presencia & 659 de 676 ráfagas con tres picos (97 \%) \\
\hline
Energía acústica en \mbox{250--1000 Hz} & 0,986 \\
\hline
\end{tabular}
\end{table}

El resultado determinante no son las frecuencias sino su \textbf{relación con el régimen de
giro}. Las tres son múltiplos enteros de la frecuencia fundamental con una desviación inferior
al 0,05 \%, y el cociente correspondiente al noveno armónico vale 9,0042 con un coeficiente de
variación del 0,137 \% sobre las 429 ráfagas en que ese armónico ocupa el pico de frecuencia
más alta.

Una resonancia estructural no produce tres componentes simultáneas en múltiplos enteros
exactos de una frecuencia que además varía. El argumento no se apoya en la estabilidad de una
sola componente, que sería discutible, sino en la \textbf{coincidencia de tres} con la serie
armónica del giro. La familia observada la genera por tanto el propio giro del motor.

Procede consignar una precisión sobre el número de componentes. El firmware publica tres picos
espectrales por eje, uno de los cuales corresponde a la fundamental, de modo que en cada
ráfaga solo caben dos armónicos. La pareja concreta que se registra varía: el octavo y el
noveno en 401 ráfagas, el noveno y el décimo en 235. El fenómeno es estable y lo que varía es
cuál de los tres armónicos queda fuera del espacio disponible. Esta circunstancia sitúa el
límite de la instrumentación empleada: con tres picos no se puede caracterizar una familia
armónica de más de dos componentes.

\section{Confirmación por modalidades independientes}

La disponibilidad de un activo de referencia permite descartar que la componente sea un
artefacto de la cadena de medida. El cuadro~\ref{table:contraste-modalidades} contrasta ambos
activos.

\begin{table}[ht!]
\caption{Contraste entre el activo con fallo y el de referencia, con el mismo firmware.}
\label{table:contraste-modalidades}
\centering
\begin{tabular}{|p{5.6cm}|p{3.2cm}|p{3.2cm}|}
\hline
\textbf{Magnitud} & \textbf{Con fallo} & \textbf{Referencia} \\
\hline
Energía acústica en \mbox{0--250 Hz} & 0,013 & 0,789 \\
\hline
Energía acústica en \mbox{250--1000 Hz} & \textbf{0,986} & 0,192 \\
\hline
Nivel acústico relativo & \textbf{1732} & 897 \\
\hline
Picos espectrales significativos & \textbf{3} & \textbf{1} \\
\hline
Ráfagas con componente en \mbox{380--460 Hz} & 627/676 (93 \%) & 0/45 (0 \%) \\
\hline
\end{tabular}
\end{table}

El reparto de energía acústica resulta \textit{invertido} entre ambos activos, y las
componentes que detecta el acelerómetro caen precisamente en la banda donde se concentra el
98,6 \% de la energía acústica del activo con fallo. Las dos modalidades coinciden sin compartir sensor ni
cadena de acondicionamiento, lo que constituye la justificación experimental del planteamiento
multimodal expuesto en el capítulo~\ref{cap:estado-arte}: una única modalidad habría aportado
un indicio, mientras que la coincidencia de dos independientes lo respalda.

\section{Consecuencias de diseño y limitaciones del contraste}

Procede consignar tres circunstancias que este resultado revela sobre el propio diseño del
sistema, por su valor metodológico.

La primera es que la componente \textbf{no aparece} en el valor eficaz, el valor de pico ni la
kurtosis, dado que el filtro paso bajo de \mbox{150 Hz} descrito en el
capítulo~\ref{cap:diseno} la elimina de los estadísticos temporales. Su detección depende por
completo de los picos espectrales, que se calculan sobre la señal sin filtrar. La decisión de
no filtrar el espectro, adoptada para poder caracterizar un artefacto que entonces se
atribuía al acoplamiento, resultó ser la que hizo posible detectar el fallo.

La segunda refuerza la anterior. Se valoró como alternativa acotar la búsqueda del pico
espectral a la banda considerada fiable del montaje, del orden de \mbox{150 Hz}, y se
descartó por no depender de acertar con un límite fijo. Esa alternativa habría situado la
componente del fallo fuera de la banda examinada y el fenómeno no se habría registrado.

La tercera afecta a la interpretación del artefacto. La componente se atribuyó inicialmente a
una resonancia del acoplamiento del sensor, hipótesis coherente con las observaciones
disponibles en ese momento: aparecía en distintas posiciones del sensor y su frecuencia se
desplazaba ligeramente. El indicio que llevó a revisarla fue la información del operador sobre
el tono audible, ajena a los datos. Se documenta porque ilustra un límite del análisis basado
exclusivamente en los registros: la prueba de enganche al régimen de giro, que resuelve la
cuestión de forma concluyente, no se planteó hasta que hubo un motivo externo para dudar de la
interpretación establecida.

El activo con fallo y el de referencia son máquinas distintas, de modo que el contraste no
permite separar el efecto del fallo del de la variabilidad entre ejemplares. La magnitud de la
diferencia observada, un reparto de energía acústica de 0,03 frente a 0,99, resulta difícil de
atribuir a esa variabilidad, pero la afirmación rigurosa exige la comparación dentro de un
mismo activo que proporcionan las campañas de fallo inducido previstas.

Debe consignarse asimismo que el modo de fallo no está identificado. Se ha caracterizado su
firma —una familia de armónicos del giro, presente en vibración y en acústica— sin que ello
permita, con la evidencia disponible, atribuirla a un mecanismo concreto.

\section{Caracterización del comportamiento nominal}

La campaña de referencia comprende \mbox{19,04 h} continuas sobre el activo en uso doméstico
normal, con 1868 ráfagas de las que 467 resultan utilizables tras descartar las capturadas con
el compresor detenido y las que no superan el filtro de calidad. El
cuadro~\ref{table:firma-nominal} recoge la firma operativa resultante.

\begin{table}[ht!]
\caption{Firma operativa del activo en estado nominal (campaña EXP-005, 467 ráfagas en 22 episodios).}
\label{table:firma-nominal}
\centering
\begin{tabular}{|p{5.4cm}|r|r|}
\hline
\textbf{Magnitud} & \textbf{Mediana} & \textbf{CV} \\
\hline
Frecuencia de giro & \mbox{49,16 Hz} & 4,2 \% \\
\hline
Valor eficaz, banda \mbox{0--150 Hz} & \mbox{2,00 m/s$^2$} & 35,1 \% \\
\hline
Valor de pico & \mbox{2,97 m/s$^2$} & 33,3 \% \\
\hline
Kurtosis & 1,518 & 16,1 \% \\
\hline
Factor de cresta & 1,470 & 11,0 \% \\
\hline
Picos espectrales significativos & 1 & — \\
\hline
Energía acústica \mbox{0--250 Hz} & 0,866 & 14,0 \% \\
\hline
\end{tabular}
\end{table}

Interesa el contraste entre los coeficientes de variación. Las magnitudes con unidades —valor
eficaz, valor de pico y amplitud dominante— presentan dispersiones en torno al 35 \%, mientras
que las adimensionales se sitúan entre el 4 \% y el 16 \%. La diferencia no es casual: las
primeras dependen del punto del ciclo térmico en que se capture la ráfaga y de la carga del
compresor, mientras que las segundas describen la \textit{forma} de la señal. Es una
justificación experimental adicional de la elección de características expuesta más adelante.

Antes de interpretar la distribución horaria procede una advertencia. El nodo carece de reloj de
tiempo real y las marcas de tiempo las asigna el concentrador al recibir cada mensaje, cuyo reloj
no está sincronizado con una referencia externa. Los \textbf{intervalos y las duraciones son
válidos}, dado que proceden de un único reloj, pero la \textbf{alineación absoluta con la hora
civil no está verificada}. Toda afirmación que dependa de la hora del día debe leerse con ese
margen, y en particular la atribución de un tramo a la noche.

El comportamiento cíclico se aprecia con nitidez en el tramo que las marcas sitúan entre las
00:00 y las 10:00, cuando cesa el uso del aparato: nueve arranques de \mbox{30 min} de duración mediana
separados por \mbox{83,8 min}, lo que sitúa el ciclo de trabajo en torno al 36 \%. El tramo de
tarde y noche presenta en cambio arranques irregulares de 4 a \mbox{9 min}, atribuibles a las
aperturas de la puerta. Esa heterogeneidad no es un defecto de la campaña sino la variabilidad
que el modelo debe tolerar, y es precisamente lo que la validación cruzada por episodios pone a
prueba.

El canal lento aporta 65 636 tramas sobre el mismo intervalo. La temperatura ambiente osciló
entre \mbox{22,4 °C} y \mbox{34,4 °C} y la del motor entre \mbox{27,5 °C} y \mbox{54,6 °C}, con
un diferencial medio de \mbox{11,69 °C}. Debe recordarse que esa temperatura la proporciona el
sensor integrado del acelerómetro y no una sonda sobre la carcasa del motor, de modo que
constituye un indicador de proximidad y no una medida del devanado; su utilidad práctica reside
en distinguir de forma inequívoca el compresor en marcha del detenido.

\section{Detección de anomalías}

Se dispone de un resultado preliminar, obtenido con el conjunto de referencia todavía en
curso de captura. Se presenta con las reservas que su tamaño de muestra impone, porque la
conclusión metodológica que de él se extrae no depende de esas reservas.

Un primer ajuste sobre las características con unidades físicas —valor eficaz, valor de pico
y kurtosis— identificó como anómalas el 99,8 \% de las ráfagas del activo con fallo. La cifra
\textbf{no mide capacidad diagnóstica}. El nivel de vibración de los dos activos difiere en un
factor 5,5, de modo que cualquier característica dimensional separa las dos máquinas antes de
separar el estado de una de ellas; el detector alcanzaba ese resultado sin haber empleado
ninguna de las características que contienen la firma del fallo. Se consigna porque ilustra un
modo de error al que este tipo de contraste está expuesto: una métrica excelente obtenida
sobre el atributo equivocado.

La formulación adoptada emplea en su lugar \textbf{características adimensionales}, invariantes
a la escala de la máquina: la kurtosis, el factor de cresta, la relación de cada pico
espectral secundario con la fundamental, su amplitud relativa y el reparto de energía acústica
por bandas. El cuadro~\ref{table:separacion-caracteristicas} recoge la separación que cada una
alcanza, expresada en desviaciones típicas.

\begin{table}[ht!]
\caption{Separación entre el activo de referencia y el activo con fallo por característica adimensional, sobre 467 y 656 ráfagas.}
\label{table:separacion-caracteristicas}
\centering
\begin{tabular}{|p{5.4cm}|r|r|r|}
\hline
\textbf{Característica} & \textbf{Nominal} & \textbf{Fallo} & \textbf{Separación} \\
\hline
Número de picos significativos & 1,000 & 3,000 & \textbf{6,90 sd} \\
\hline
Energía acústica \mbox{250--1000 Hz} & 0,122 & 0,986 & 4,07 sd \\
\hline
Energía acústica \mbox{0--250 Hz} & 0,866 & 0,013 & 3,88 sd \\
\hline
Relación armónica del pico superior & 1,997 & 9,005 & 3,17 sd \\
\hline
Relación armónica del pico intermedio & 2,999 & 8,004 & 2,21 sd \\
\hline
Amplitud relativa del pico superior & 0,015 & 0,664 & 1,96 sd \\
\hline
Factor de cresta & 1,470 & 1,975 & 0,94 sd \\
\hline
Kurtosis & 1,518 & 1,697 & 0,57 sd \\
\hline
\end{tabular}
\end{table}

La selección del modelo se resolvió comparando siete candidatos en el mismo punto de operación
—umbral en el cuantil 0,05 de las puntuaciones de entrenamiento— con arreglo a cinco criterios
aplicados en orden: comportamiento sobre condiciones de operación no vistas, tasa de falsos
positivos, tasa de detección, número de hiperparámetros y viabilidad de embarcar la inferencia
en el microcontrolador.

Debe señalarse en primer lugar que \textbf{la tasa de detección no discrimina}: los siete
candidatos la sitúan en el 100 \%. Con un factor 12,5 de diferencia de nivel de vibración entre
los dos activos, detectar este fallo es sencillo y esa cifra no mide capacidad diagnóstica.
Presentarla como logro sería el mismo error que el detector dimensional descrito más arriba.

\section{Validación cruzada por episodios}

La validación se realizó dejando fuera un \textbf{episodio de marcha completo}, ajustando el
modelo sobre los restantes y midiendo los falsos positivos sobre una condición de operación que
el modelo no había visto. Es la única forma de estimar si el detector generaliza entre arranques
del compresor y no solamente entre ráfagas correlacionadas del mismo arranque, y exigió la
campaña de referencia prolongada: con un único episodio no era posible.

El cuadro~\ref{table:validacion-episodios} recoge el resultado sobre los 22 episodios
disponibles.

\begin{table}[ht!]
\caption{Validación cruzada por episodios sobre el conjunto de referencia (22 episodios, 467 ráfagas).}
\label{table:validacion-episodios}
\centering
\begin{tabular}{|p{4.6cm}|r|r|r|}
\hline
\textbf{Modelo} & \textbf{Falsos pos.} & \textbf{Peor episodio} & \textbf{Detección} \\
\hline
\textbf{Regla sobre n.\textsuperscript{o} de picos} & \textbf{0,5 \%} & \textbf{8,3 \%} & 100 \% \\
\hline
Bosque de aislamiento & 3,3 \% & 33,8 \% & 100 \% \\
\hline
Regla sobre relación armónica & 4,6 \% & 18,8 \% & 96,6 \% \\
\hline
Envolvente robusta & 5,9 \% & 33,3 \% & 100 \% \\
\hline
Factor local de atipicidad & 6,2 \% & 31,1 \% & 100 \% \\
\hline
Máquina de vectores soporte & 6,7 \% & 25,0 \% & 100 \% \\
\hline
\end{tabular}
\end{table}

La columna determinante es la del peor episodio, y no la media. Una media baja con un peor caso
elevado significa que el detector falla de forma \textit{concentrada}: en operación se traduce
en una tanda de alarmas falsas consecutivas sobre un mismo arranque y no en un falso positivo
aislado cada cierto tiempo, que es una diferencia práctica considerable para quien atiende el
sistema. La regla adoptada resulta de tres a cuatro veces mejor que cualquier alternativa en esa
columna.

El modelo adoptado es por tanto una \textbf{regla sobre el número de picos espectrales
significativos}, entendiendo por significativo el que alcanza el 20 \% de la amplitud del mayor.
Vale 1 en 465 de las 467 ráfagas del activo de referencia y 3 en 656 de las 656 del activo con
fallo. El resultado se mantiene para cualquier umbral de significación entre el 15 \% y el
30 \%, y se ejecuta en el nodo con una comparación.

Procede consignar que la máquina de vectores soporte pasó de un 46 \% de falsos positivos con
el conjunto reducido inicial a un 6,3 \% con el conjunto ampliado. El diagnóstico previo era
correcto en su naturaleza: sus hiperparámetros no resultaban ajustables con 45 observaciones, y
la limitación estaba en el tamaño de muestra y no en el modelo.

\section{El filtro de calidad forma parte del detector}

El resultado de mayor consecuencia práctica de la campaña de referencia no afecta al modelo sino
a la validez de la medida, y no era observable con las campañas breves anteriores.

\textbf{Los reintentos del bus de comunicación con el acelerómetro reproducen la firma del
fallo sobre un activo sano.} El cuadro~\ref{table:reintentos-artefacto} lo recoge sobre el
activo de referencia, cuyo estado es nominal.

\begin{table}[ht!]
\caption{Efecto de los reintentos del bus sobre las características, en el activo nominal.}
\label{table:reintentos-artefacto}
\centering
\begin{tabular}{|r|r|r|r|}
\hline
\textbf{Reintentos} & \textbf{Ráfagas} & \textbf{Con más de un pico} & \textbf{Fundamental} \\
\hline
0 & 261 & 0,4 \% & \mbox{49,15 Hz} \\
\hline
1--3 & 206 & 0,5 \% & \mbox{49,17 Hz} \\
\hline
4--5 & 53 & 8 \% & \mbox{49,25 Hz} \\
\hline
6--10 & 48 & 31 \% & \mbox{49,16 Hz} \\
\hline
11--20 & 46 & \textbf{93 \%} & \textbf{\mbox{20,03 Hz}} \\
\hline
más de 20 & 69 & \textbf{94 \%} & \textbf{\mbox{16,03 Hz}} \\
\hline
\end{tabular}
\end{table}

El mecanismo es que una muestra corrupta inyecta ruido de banda ancha en la ráfaga, con lo que
varios coeficientes del espectro superan el umbral de significación y la estimación de la
frecuencia fundamental deja de corresponder al régimen de giro. De ello se sigue que el filtro
de calidad no es un preproceso sino \textbf{parte del detector}: el nodo debe negarse a emitir
veredicto sobre una ráfaga cuyo número de reintentos exceda el límite establecido, en lugar de
juzgarla. El límite se calibró en tres reintentos a partir del cuadro anterior, valor que
conserva 467 ráfagas con un 0,4 \% de artefactos —la misma proporción que exigir cero
reintentos— con un 79 \% más de observaciones.

\section{La firma del fallo no es ese artefacto}

La coincidencia entre el artefacto y la firma obliga a descartar que el hallazgo de EXP-003 sea
una manifestación del mismo defecto. La comprobación consiste en restringir ambos conjuntos a
las ráfagas capturadas \textbf{sin ningún reintento}, y su resultado figura en el
cuadro~\ref{table:cero-reintentos}.

\begin{table}[ht!]
\caption{Ambos activos restringidos a ráfagas sin ningún reintento del bus.}
\label{table:cero-reintentos}
\centering
\begin{tabular}{|p{4.4cm}|r|p{3.4cm}|r|}
\hline
\textbf{Conjunto} & \textbf{Ráfagas} & \textbf{Picos: 1 / 2 / 3} & \textbf{Fundamental} \\
\hline
Activo de referencia & 261 & 260 / 0 / 1 & \mbox{49,15 Hz} \\
\hline
Activo con fallo & 443 & 0 / 13 / 430 & \mbox{49,76 Hz} \\
\hline
\end{tabular}
\end{table}

La separación se mantiene íntegra. Los cocientes armónicos del activo con fallo sobre ese
subconjunto valen \textbf{8,0016 con un coeficiente de variación del 0,029 \%} y 9,0035 con un
0,106 \%, cifras más precisas que las obtenidas sobre el conjunto completo. A ello se añade un
argumento de sentido contrario que resulta concluyente: la firma del activo con fallo es
\textit{idéntica} con reintentos y sin ellos —8,003 frente a 7,998— mientras que las
características del activo de referencia se derrumban en presencia de reintentos. Un artefacto
del bus no puede producir un valor que el propio bus no altera.


\section{Revisión de la elección de modelo}

La elección expuesta hasta aquí se revisó, y procede consignar el motivo porque afecta a una
conclusión y no a un detalle.

La primera conclusión fue que la regla sobre el número de picos superaba a los modelos de
aprendizaje automático. \textbf{Esa conclusión es engañosa.} La comparación no era equitativa:
a la regla se le proporcionaba una característica seleccionada con cuidado, mientras que a los
modelos se les proporcionaban las quince, entre ellas varias sin capacidad de separación y con
dispersión elevada en el conjunto nominal, circunstancia que penaliza a todo modelo basado en
distancias o en densidades.

Repetida la comparación proporcionando a cada modelo \textbf{únicamente esa característica},
los seis candidatos convergen al mismo 8,3 \% de falsos positivos en el peor episodio. Con una
sola dimensión existe una sola frontera que localizar, de modo que todos localizan la misma. El
mérito no corresponde al algoritmo sino a la característica.

De ello se sigue la objeción determinante. Esa característica se seleccionó \textbf{conociendo
cuál era el fallo}, de modo que su ventaja no se transfiere a un fallo distinto, y el análisis
de sensibilidad expuesto más arriba lo confirma: la regla resulta ciega a cuatro de las cinco
direcciones examinadas. Un detector de anomalías tiene por objeto precisamente los fallos que
no se conocen de antemano, y optimizar su elección contra el único fallo disponible constituye
sobreajuste al mismo.

Procede además matizar la lectura del 25 \% de falsos positivos de la envolvente robusta en el
peor episodio, cifra que se presentó de forma que induce a error. Ese porcentaje corresponde a
\textbf{tres ráfagas de doce} en uno de los episodios cortos e irregulares del tramo de tarde.
Ponderada por número de ráfagas —que es la magnitud que se observaría en operación— la tasa es
del \textbf{5,6 \%}, y once de los veintidós episodios no producen ningún falso positivo.

\textbf{El detector principal es por tanto la envolvente robusta} sobre el conjunto completo de
características, y la regla sobre el número de picos se conserva como confirmación de alta
confianza para fallos de naturaleza armónica. Las dos son ejecutables en el microcontrolador
—una forma cuadrática con matriz precalculada de 900 bytes en el primer caso, una comparación
en el segundo—, de modo que se embarcan ambas con significados distintos: la envolvente
señala que el estado de la máquina se ha alejado del de referencia; la regla, que la desviación
corresponde a una familia armónica. Ambas requieren histéresis, es decir la exigencia de varias
ráfagas anómalas consecutivas antes de notificar: el 8,3 \% de falsos positivos del peor
episodio corresponde a ráfagas aisladas, y exigir tres consecutivas lo reduce en dos órdenes de
magnitud bajo el supuesto de independencia.

Debe señalarse por último que existe una dirección de fallo que \textbf{ningún} modelo sobre el
canal de vibración puede cubrir: la obstrucción del sistema de ventilación no altera la firma
vibratoria. Su detección corresponde al canal lento, a través del diferencial térmico y de la
prolongación del ciclo de trabajo, lo que constituye una justificación adicional del
planteamiento multimodal.



\section{Publicación del veredicto y reducción del tráfico}

El veredicto se publica en un tema propio y no como campos añadidos al mensaje de ráfaga, y la
razón es de integridad del conjunto de datos: el registrador escribe la cabecera del fichero
únicamente al crearlo, de modo que añadir campos desplazaría las columnas de las filas
posteriores y dividiría la serie histórica en dos conjuntos no comparables. Se verificó que el
mensaje de ráfaga conserva sus 45 campos en el mismo orden.

El mensaje de veredicto ocupa \mbox{124 B} en el peor caso, frente a los del orden de
\mbox{1050 B} del mensaje de características. Un consumidor que únicamente requiera el estado
del activo recibe por tanto \textbf{un orden de magnitud menos de tráfico}, lo que constituye la
materialización del objetivo enunciado en el capítulo~\ref{cap:objetivos}.

El veredicto contempla \textbf{tres} estados y no dos. El tercero declara que la medida no es
apta para decidir, y se emite con más reintentos del bus que el límite establecido o con el
compresor detenido. No constituye un estado de salud intermedio: su función es que el nodo se
niegue a juzgar antes que juzgar sobre una medida que, según se expuso más arriba, reproduce
artificialmente la firma del fallo.

\section{Efecto de la histéresis sobre la carga de alarmas}

Un detector que notifica ante una única ráfaga anómala resulta inservible en operación. El
cuadro~\ref{table:histeresis} recoge el número de avisos que se emitirían sobre la secuencia
completa de ráfagas según el número de ráfagas anómalas consecutivas exigidas.

\begin{table}[ht!]
\caption{Avisos emitidos sobre las 1161 ráfagas según la histéresis exigida.}
\label{table:histeresis}
\centering
\begin{tabular}{|r|r|r|}
\hline
\textbf{Ráfagas exigidas} & \textbf{Avisos falsos} & \textbf{Avisos sobre el fallo} \\
\hline
1 & 22 & 1 \\
\hline
2 & 4 & 1 \\
\hline
\textbf{3} & \textbf{0} & \textbf{1} \\
\hline
4 & 0 & 1 \\
\hline
\end{tabular}
\end{table}

Exigiendo tres ráfagas consecutivas se emiten \textbf{cero avisos falsos} sobre las 505 ráfagas
del activo de referencia, y el fallo se notifica. El 5,1 \% de ráfagas marcadas como anómalas no
produce ningún aviso porque se trata de observaciones aisladas.

De ello se sigue una precisión metodológica: la magnitud que estima la carga de alarmas en
operación es el número de \textit{avisos emitidos} y no la fracción de ráfagas marcadas. Ambas
difieren en dos órdenes de magnitud.

\section{Respuesta ante una perturbación del entorno: apertura de puerta}

Se ejecutó una campaña abriendo la puerta del aparato entre las \mbox{12:09} y las
\mbox{12:11}, con el activo de referencia en estado nominal. Procede advertir de entrada que
esta campaña \textbf{no se prerregistró}, a diferencia de la de obstrucción de ventilación, de
modo que sus criterios de evaluación no se declararon con antelación.

El cuadro~\ref{table:puerta} recoge la secuencia registrada en el tema de estado.

\begin{table}[ht!]
\caption{Respuesta del detector embarcado ante la apertura de puerta. Umbral del nodo: $-1{,}4395$.}
\label{table:puerta}
\centering
\begin{tabular}{|l|c|r|c|c|}
\hline
\textbf{Instante} & \textbf{Estado} & \textbf{Puntuación} & \textbf{Racha} & \textbf{Aviso} \\
\hline
12:08:40 & \texttt{nominal} & $-1{,}118$ & 0 & 0 \\
12:09:10 & \texttt{nominal} & $-1{,}170$ & 0 & 0 \\
\hline
12:09:40 & \textbf{\texttt{anomaly}} & $\mathbf{-2{,}156}$ & 1 & 0 \\
12:10:10 & \textbf{\texttt{anomaly}} & $\mathbf{-2{,}100}$ & 2 & 0 \\
12:10:40 & \textbf{\texttt{anomaly}} & $\mathbf{-1{,}862}$ & \textbf{3} & \textbf{1} \\
\hline
12:11:10 & \texttt{nominal} & $-1{,}278$ & 0 & 0 \\
12:11:40 & \texttt{nominal} & $-1{,}048$ & 0 & 0 \\
\hline
\end{tabular}
\end{table}

La magnitud del desplazamiento de la puntuación invita a atribuirlo a un efecto térmico. El
examen de las características individuales lo desmiente, y el
cuadro~\ref{table:puerta-features} lo recoge.

\begin{table}[ht!]
\caption{Desplazamiento de cada característica durante las tres ráfagas anómalas, en desviaciones típicas del conjunto de referencia.}
\label{table:puerta-features}
\centering
\begin{tabular}{|p{6.6cm}|r|}
\hline
\textbf{Característica} & \textbf{Desplazamiento} \\
\hline
Energía acústica \mbox{250--1000 Hz} & \textbf{+2,70 sd} \\
\hline
Energía acústica \mbox{0--250 Hz} & \textbf{$-$2,64 sd} \\
\hline
Amplitud dominante relativa & +1,30 sd \\
\hline
Valor eficaz relativo & +1,25 sd \\
\hline
Diferencial térmico relativo & +1,23 sd \\
\hline
Kurtosis, factor de cresta, relaciones armónicas, n.\textsuperscript{o} de picos & sin desplazamiento apreciable \\
\hline
\end{tabular}
\end{table}

Lo que el nodo detectó fue un \textbf{cambio del entorno acústico}: la energía sonora se
desplazó de la banda inferior a la intermedia, dado que el micrófono registró el recinto con la
puerta abierta. La \textbf{firma vibratoria no se alteró}, y es la firma vibratoria la que
describe el estado mecánico del activo.

El activo se encontraba sano y la perturbación la introdujo la apertura de una puerta, suceso
que en un aparato de uso doméstico es el más frecuente de su vida operativa. En explotación,
este episodio sería un \textbf{falso positivo} y no una alarma útil.

\textbf{El aviso inmediato no aporta por tanto validación dentro de una misma máquina.} Abrir
una puerta no degrada el compresor. Procede matizar, no obstante, que la campaña sí registró
sobre el propio activo de referencia una \textit{condición} anómala real —la sobrecarga térmica
que se expone en el apartado siguiente— si bien se trata de una sobrecarga operativa y no de un
fallo mecánico. La limitación de mayor peso del trabajo permanece en consecuencia intacta en lo
que respecta a los \textbf{fallos}: el contraste entre el activo averiado y el de referencia
sigue estableciéndose entre máquinas distintas, y solo un fallo inducido sobre el activo de
referencia lo resuelve.

Tampoco cabe afirmar una alteración del régimen de ciclado. El compresor había arrancado a las
\mbox{11:32}, treinta y siete minutos antes de la apertura, de modo que la marcha continua no es
atribuible a la perturbación; y la captura concluye a las \mbox{12:41}, con lo que treinta
minutos de observación posterior no permiten establecer un cambio en un ciclo cuyo periodo es de
unos \mbox{85 min}.

\section{El efecto térmico sí fue real, y el detector lo señaló con retraso}

Procede matizar lo anterior, porque la perturbación tuvo una consecuencia física sostenida que
conviene separar del transitorio acústico. El cuadro~\ref{table:puerta-termico} recoge la
evolución del estado térmico.

\begin{table}[ht!]
\caption{Evolución térmica y respuesta del detector alrededor de la apertura.}
\label{table:puerta-termico}
\centering
\begin{tabular}{|p{4.6cm}|r|r|r|}
\hline
\textbf{Ventana} & \textbf{Diferencial} & \textbf{Temp. motor} & \textbf{Ráfagas marcadas} \\
\hline
Marcha previa (11:32--12:09) & \mbox{15,26 °C} & \mbox{41,0 °C} & 4 de 18 \\
\hline
Apertura (12:09--12:11) & \mbox{19,48 °C} & \mbox{46,7 °C} & 3 de 3 \\
\hline
De +0 a +30 min & \mbox{20,64 °C} & \mbox{48,8 °C} & 10 de 45 \\
\hline
De +30 a +47 min & \mbox{23,81 °C} & \mbox{50,3 °C} & \textbf{23 de 23} \\
\hline
\end{tabular}
\end{table}

El diferencial térmico subió de \mbox{15,26 °C} a \mbox{23,81 °C} y la temperatura del motor de
\mbox{41,0 °C} a \mbox{50,3 °C}, de forma monótona y creciente durante los cuarenta y siete
minutos posteriores: el efecto es real y es la consecuencia física esperable de la perturbación.

Y el detector lo señaló con una progresión que la tabla hace evidente. En la primera media hora
marcó 10 de 45 ráfagas; en la siguiente, con el diferencial ya \mbox{8,5 °C} por encima de la
referencia, marcó \textbf{las 23 de 23}, y la racha consecutiva alcanzó 25 al término de la
captura. No es una reacción esporádica sino una \textbf{detección sostenida}.

De ello se sigue la lectura del episodio, que contempla \textbf{dos fenómenos distintos con
significado opuesto}.

El primero es el aviso de las \mbox{12:10:40}, a los noventa segundos de la apertura. La prueba
de ablación lo atribuye al transitorio acústico, y constituye un \textbf{falso positivo} sobre
un suceso operativo normal: en un aparato de uso doméstico, abrir la puerta es el suceso más
frecuente de su vida operativa.

El segundo es la detección sostenida de la sobrecarga térmica, que se manifiesta con
\textbf{treinta minutos de retraso} respecto de la perturbación y de forma continua a partir de
entonces. Ese retraso no es un defecto sino la constante de tiempo del fenómeno: el estado
térmico de un compresor no cambia en noventa segundos. Y sí constituye la detección de una
condición anómala real del propio activo de referencia, si bien de una \textbf{sobrecarga
operativa} y no de un fallo mecánico.

Procede consignar dos reservas. La captura concluye a las \mbox{12:58} con el detector aún en
estado de anomalía y el diferencial todavía creciendo, de modo que no se ha observado el retorno
al estado nominal, que es la mitad del ciclo que una campaña completa debería recoger. Y el
activo ya presentaba un diferencial elevado antes de la apertura por llevar treinta y siete
minutos de marcha continua, circunstancia que impide separar limpiamente el efecto de la
perturbación del de la marcha prolongada previa.

\section{El coste del canal acústico}

Dos resultados de esta campaña son sólidos y ambos afectan al diseño. El primero es que
\textbf{el mecanismo de histéresis se comporta según su diseño}: se acumularon exactamente tres
ráfagas anómalas consecutivas y el aviso se emitió en la tercera y no antes, según se aprecia en
el cuadro~\ref{table:puerta}. El segundo es la cuantificación del coste del canal acústico. Anulando las tres bandas y
recalculando la puntuación sobre las ráfagas de la captura cerrada a las \mbox{12:58} del 28 de
agosto, la tasa de ráfagas
marcadas desciende del \textbf{14,0 \% al 5,7 \%}, y el episodio de la apertura deja de
detectarse por completo: ninguna de sus tres ráfagas resulta marcada. El canal acústico aporta
en consecuencia \textbf{8,3 puntos} de tasa de ráfagas marcadas y es el único responsable de la
detección de la apertura.

Debe matizarse el mecanismo. Únicamente el 10 \% de las ráfagas anómalas presenta un valor
extremo de energía en la banda intermedia, frente al 5 \% que correspondería por construcción, de
modo que la contribución del canal no procede mayoritariamente de sucesos intensos como el de la
apertura sino de una dispersión distribuida que la campaña de referencia no capturó.

El canal acústico presenta por tanto un \textbf{compromiso con las dos caras medidas}: aportó la
confirmación independiente del fallo real, procedente de un sensor que no comparte cadena con el
acelerómetro, y cuesta 8,3 puntos de tasa de marcado frente a perturbaciones del entorno. La
elección entre ambas no es de naturaleza técnica —depende de si el activo opera en un entorno
acústicamente estable— y queda declarada y no resuelta. Los datos sugieren una alternativa sin
ensayar: exigir que la desviación se manifieste \textit{también} en el canal de vibración antes
de notificar, con lo que el canal acústico conservaría su función confirmatoria sin poder
disparar por sí solo.

\section{Procedimientos de verificación y selección}

Tres procedimientos sostienen los resultados anteriores y se detallan en el
anexo~\ref{anexo:verificacion} por su extensión: el protocolo de selección de modelo con
separación estricta entre desarrollo y prueba, el análisis de sensibilidad frente a modos de
fallo no observados, y la verificación del código embarcado frente a la implementación de
referencia.

De los tres conviene retener aquí lo que condiciona la interpretación de los resultados.

El \textbf{protocolo de selección} se estableció tras advertir que el procedimiento inicial
había empleado el conjunto de evaluación para adoptar decisiones de diseño, lo que constituye
sesgo de espionaje de datos. Con separación estricta y partición cronológica por episodios, el
modelo seleccionado alcanza un \mbox{7,8 \%} de falsos positivos sobre episodios no vistos y un
\mbox{100 \%} de detección. El protocolo \textbf{selecciona un modelo distinto} del elegido por
el procedimiento anterior, discrepancia que constituye la evidencia de que aquella elección
estaba contaminada.

El \textbf{análisis de sensibilidad} indica que el modelo reaccionaría a las cinco direcciones de
fallo examinadas, si bien emplea perturbaciones sintéticas del conjunto nominal y \textbf{no
constituye evidencia experimental}. Su valor es localizar puntos ciegos, y el principal que
localiza es que una regla sobre una única característica resulta ciega a cuatro de las cinco.

La \textbf{verificación del código embarcado} se realizó sobre las 1161 ráfagas reales de las dos
campañas, con \textbf{cero veredictos discrepantes} frente a la implementación de referencia. De
ella se sigue que la inferencia no requiere entorno de ejecución alguno: los entornos disponibles
para microcontroladores operan sobre redes neuronales, y ninguno de los modelos considerados
pertenece a esa familia.

\section{Alcance de la evidencia obtenida}

Procede acotar el tamaño de muestra efectivo, que no coincide con el número de ráfagas. Las
ráfagas se publican cada \mbox{30 s} y se agrupan en episodios de marcha, de modo que las de un
mismo episodio describen la misma condición de operación y están fuertemente correlacionadas. El
conjunto de referencia procede de \textbf{22 episodios} sobre \mbox{19,04 h}, lo que hace
posible la validación cruzada expuesta más arriba. El del activo con fallo procede en cambio de
\textbf{dos} episodios aprovechables, uno de \mbox{30 min} y otro de \mbox{367 min}: sus 656
ráfagas no constituyen 656 observaciones independientes, y ese desequilibrio no se corrige
alargando la captura, dado que ese activo permaneció en marcha el 95 \% del tiempo frente al
37 \% del activo de referencia. Obtener episodios adicionales exigiría provocar una parada.

Tres limitaciones acotan lo que la evidencia permite afirmar.

La primera es que \textbf{la validación por episodios se aplica al lado nominal y no al del
fallo}, por el desequilibrio de episodios ya expuesto.

La segunda es que \textbf{en el activo secundario el modo de fallo exacto no está clasificado},
disponiéndose de una firma armónica caracterizada y no de un diagnóstico desagregado.

La tercera afecta al rendimiento de la captura: del conjunto de referencia solo resulta
utilizable el 25 \% de las ráfagas. La causa principal es el ciclo de trabajo —el compresor está
en marcha el 37 \% del tiempo y una ráfaga con el compresor detenido no contiene vibración que
analizar— y la secundaria son los reintentos del bus, que descartan un 32 \% adicional de las
ráfagas en marcha. Una hora de captura rinde por tanto unas 25 ráfagas utilizables, cifra sobre
la que debe dimensionarse toda campaña y no sobre la cadencia de publicación.

\section{Objetivos no alcanzados}

Procede enunciar de forma explícita los objetivos que no se han alcanzado, con su
justificación.

\textbf{La clasificación automática de los estados de salud no se ha implementado como tal.} La
propuesta de trabajo la enuncia entre las tareas del motor de análisis, y lo obtenido es un
detector de \textit{una clase} que distingue el estado nominal de cualquier desviación respecto
de él, sin atribuirla a un estado concreto. La razón es de disponibilidad de datos y no de
tiempo: un clasificador exige ejemplos etiquetados de cada estado que deba distinguir, y se
dispone de un catálogo limitado de modos de fallo. Con esa evidencia, un clasificador multiclase
no sería ajustable ni evaluable de forma rigurosa.

Debe añadirse que la detección de una clase es el planteamiento que la literatura prescribe
precisamente para esta situación, según se expuso en el capítulo~\ref{cap:estado-arte}: los
fallos son escasos, variados y no se conocen de antemano. La distinción entre estado nominal y
estado anómalo sí queda resuelta, de modo que el objetivo se alcanza en su formulación binaria y
no en la multiclase.

\textbf{El aprendizaje y ajuste de los parámetros del modelo se realiza en el computador y no en el propio nodo.} Si bien la inferencia y el diagnóstico se ejecutan íntegramente en el microcontrolador en tiempo real, la estimación inicial de medianas y el cálculo del modelo LOF de referencia se exportan como cabecera C++ desde el entorno de análisis tras una fase de calibración previa. Queda como trabajo futuro dotar al nodo de una rutina de autocalibración autónoma en el arranque.

\textbf{No se ha determinado la resonancia del acoplamiento del sensor}, cuya estimación inicial
resultó ser en realidad la firma de un fallo. El corte del filtro paso bajo en \mbox{150 Hz}
conserva por ello el carácter de cota conservadora sin verificación experimental.

Finalmente, \textbf{la caracterización de la firma armónica del fallo es incompleta} por una
limitación de la instrumentación que el propio análisis puso de manifiesto: con tres picos
espectrales por eje, uno de los cuales corresponde a la fundamental, solo caben dos armónicos
por ráfaga, mientras que la familia observada consta de tres.
